\chapter{Unconstrained Model Predictive Control}
\label{chap:problem8}

This chapter designs the unconstrained MPC for the linear four-tank model of
Problem~5 and applies the same design function to the identified model of
Problem~4. Without input or output constraints the problem is an unconstrained
convex quadratic program, so the design reduces to the stacked prediction
matrices and the closed form $U^\star=-H^{-1}g$, which the controller obtains
through the quadratic-programming interface of Problem~7. Two designs are
simulated, one for each estimator of Problem~6: a four-state filter and a
six-state bias-augmented filter.

\section{Model and Design Parameters}
\label{sec:p8-model}

All signals are deviations from the operating point
\begin{equation}
 u_0=\begin{bmatrix}300\\300\end{bmatrix}\ \mathrm{cm^3/s},
 \qquad
 d_0=\begin{bmatrix}20\\20\end{bmatrix}\ \mathrm{cm^3/s},
 \qquad
 h_s=\begin{bmatrix}28.4621\\41.4578\\4.0950\\6.6333\end{bmatrix}\ \mathrm{cm},
 \qquad
 T_s=5\ \mathrm{s},
 \label{eq:p8-operating-point}
\end{equation}
with $h_s$ the steady-state level. The state is the tank mass deviation
$x_k=m_k-m_s$, the input is $u_k=F_k-u_0$, the unmeasured inflow is
$d_k=D_k-d_0$, and $y_k$ and $z_k$ are the measured and controlled lower-tank
level deviations. The model is
\begin{align}
 x_{k+1} &= A x_k + B u_k + G d_k, \label{eq:p8-model-state}\\
 y_k &= C x_k + v_k, \label{eq:p8-model-meas}\\
 z_k &= C x_k, \label{eq:p8-model-out}
\end{align}
with $G=E$, zero feedthrough, and $A$, $B$, $E$ the zero-order-hold
matrices at $T_s=5$~s that were used for the Problem~6 filters.
The inflow is a step plus a zero-mean white fluctuation,
$d_k=d_k^{\mathrm{step}}+\eta_k$, and the measurement noise is $v_k$; the two
sequences are independent and Gaussian,
\begin{equation}
 \mathbb{E}[\eta_k\eta_k^T]=\sigma_F^2 I_2,
 \qquad
 \mathbb{E}[v_kv_k^T]=\sigma_V^2 I_2,
 \qquad
 \sigma_F=5\ \mathrm{cm^3/s},\quad \sigma_V=0.5\ \mathrm{cm}.
 \label{eq:p8-noise}
\end{equation}
The four-state design models $d_k$ as zero-mean white noise and carries no
disturbance state; the six-state design adds two bias states $w_k$ that absorb
the step. The design parameters are listed in Table~\ref{tab:p8-design}.

\begin{table}[htbp]
 \centering
 \caption{MPC design parameters used in the reported run. The horizon
 $N=20$ corresponds to $NT_s=100$~s.}
 \label{tab:p8-design}
 \begin{tabular}{@{}lll@{}}
 \hline
 Parameter & Value & Meaning\\
 \hline
 $T_s$ & $5$~s & sample time\\
 $N$ & $20$ & prediction horizon\\
 $W_z$ & $10 I_2$ & output weight\\
 $W_u$ & $0.01 I_2$ & input weight\\
 $W_{\Delta u}$ & $I_2$ & input-move weight\\
 $u_0$ & $[300,300]^T$~cm$^3$/s & nominal pump flow\\
 $d_0$ & $[20,20]^T$~cm$^3$/s & nominal disturbance flow\\
 \hline
 \end{tabular}
\end{table}

\section{Stacked Prediction}
\label{sec:p8-prediction}

Let $\hat x_k$ and $\hat w_k$ be the estimates after the measurement $y_k$.
Repeated substitution over the horizon gives
\begin{equation}
 x_{k+j|k}=A^j\hat x_k
 +\sum_{i=1}^{j}A^{j-i}B\,u_{k+i-1}
 +\sum_{i=1}^{j}A^{j-i}G\,\hat w_k,
 \label{eq:p8-rollout}
\end{equation}
where the disturbance estimate is held constant. Stacking the predicted outputs
gives
\begin{equation}
 Z_k=\Phi_x\hat x_k+\Phi_w\hat w_k+\Gamma U_k,
 \qquad
 Z_k=\begin{bmatrix}z_{k+1|k}\\ \vdots\\ z_{k+N|k}\end{bmatrix},
 \quad
 U_k=\begin{bmatrix}u_k\\ \vdots\\ u_{k+N-1}\end{bmatrix},
 \label{eq:p8-prediction}
\end{equation}
with the block lower triangular input-to-output matrix
\begin{equation}
 \Gamma=
 \begin{bmatrix}
  CB & 0 & \cdots & 0\\
  CAB & CB & \ddots & \vdots\\
  \vdots & \vdots & \ddots & 0\\
  CA^{N-1}B & CA^{N-2}B & \cdots & CB
 \end{bmatrix},
 \qquad
 [\Gamma]_{ji}=CA^{j-i}B,
 \label{eq:p8-gamma}
\end{equation}
and the free-response matrices
\begin{equation}
 \Phi_x=\begin{bmatrix}CA\\ CA^2\\ \vdots\\ CA^N\end{bmatrix},
 \qquad
 \Phi_w=
 \begin{bmatrix}CG\\ C(I+A)G\\ \vdots\\ C(I+\cdots+A^{N-1})G\end{bmatrix}.
 \label{eq:p8-phix-phiw}
\end{equation}
The form of $\Phi_w$ holds $\hat w_k$ constant over the horizon, which is the
prediction for the bias states of the augmented filter, whose future increments
have zero mean. When the estimator has no disturbance state, $n_w=0$ and the
disturbance term is absent. The code also supports an impulse form,
$[\Phi_w]_j=CA^{j-1}G$, for a white-disturbance model; neither design here
uses it. The reference is held at its current value,
$R_k=\mathbf{1}_N\otimes r_k$.

\section{Objective and Closed-Form Solution}
\label{sec:p8-objective}

The objective penalizes output tracking, input magnitude and input movement,
\begin{align}
 \phi_z&=\frac12\sum_{j=1}^{N}
  (z_{k+j}-r_{k+j})^T W_z^T W_z (z_{k+j}-r_{k+j}),
 \label{eq:p8-phiz}\\
 \phi_u&=\frac12\sum_{j=0}^{N-1}
  (u_{k+j}-u_{ss})^T W_u^T W_u (u_{k+j}-u_{ss}),
 \label{eq:p8-phiu}\\
 \phi_{\Delta u}&=\frac12\sum_{j=0}^{N-1}
  \Delta u_{k+j}^T W_{\Delta u}^T W_{\Delta u}\Delta u_{k+j},
 \qquad
 \Delta u_{k+j}=u_{k+j}-u_{k+j-1},
 \label{eq:p8-phidu}
\end{align}
with $u_{k-1}$ known. Stacking the weights, $\bar W_\bullet=I_N\otimes W_\bullet$,
and writing the moves as $\Delta U_k=\Lambda U_k-I_0u_{k-1}$ with
\begin{equation}
 \Lambda=
 \begin{bmatrix}
  I_2 & 0 & \cdots & 0\\
  -I_2 & I_2 & \ddots & \vdots\\
  0 & \ddots & \ddots & 0\\
  \vdots & & -I_2 & I_2
 \end{bmatrix},
 \qquad
 I_0=\begin{bmatrix}I_2\\ 0_2\\ \vdots\\ 0_2\end{bmatrix},
 \label{eq:p8-lambda-i0}
\end{equation}
the prediction error is $Z_k-R_k=\Gamma U_k-c_k$ with
$c_k=R_k-\Phi_x\hat x_k-\Phi_w\hat w_k$, and $\phi=\frac12U_k^THU_k+g_k^TU_k+\rho_k$,
where $\rho_k$ collects the terms independent of $U_k$ and
\begin{align}
 H&=\Gamma^T\bar W_z^T\bar W_z\Gamma
   +\bar W_u^T\bar W_u
   +\Lambda^T\bar W_{\Delta u}^T\bar W_{\Delta u}\Lambda,
 \label{eq:p8-H}\\
 g_k&=M_zc_k+M_u\bar U_k+M_{\Delta u}u_{k-1},
 \label{eq:p8-g}
\end{align}
where $\bar U_k=\mathbf{1}_N\otimes u_{ss}$ and
$M_z=-\Gamma^T\bar W_z^T\bar W_z$, $M_u=-\bar W_u^T\bar W_u$,
$M_{\Delta u}=-\Lambda^T\bar W_{\Delta u}^T\bar W_{\Delta u}I_0$. The gradient is
$HU_k+g_k$, so the unconstrained minimizer is
\begin{equation}
 U_k^\star=-H^{-1}g_k,
 \label{eq:p8-solution}
\end{equation}
of which only the first move,
$u_k=[I_{n_u}\ 0]U_k^\star$, is applied before the estimate and reference are
updated and the problem is solved again.

Because $W_u\succ0$ and $\Gamma$ has full column rank, $H\succ0$: for the mass
model $\Gamma\in\mathbb{R}^{40\times40}$ has rank $40$ and condition number
$25.1$, and the eigenvalues of $H$ lie between $0.0831$ and $3.979$.
\texttt{mpc\_constructors} checks this rank and the definiteness of $H$. Since
$H$ does not depend on $\hat x_k$, $\hat w_k$ or $r_k$, it is built once, while
$g_k$ is updated at every sample by \texttt{mpc\_step}, which assembles the
standard form of the Problem~7 interface,
\begin{equation}
 \min_x\ \frac12x^THx+g^Tx
 \quad\text{s.t.}\quad
 l\le x\le u,\qquad b_l\le Ax\le b_u,
 \label{eq:p8-qp-standard}
\end{equation}
and calls \texttt{qpsolver(H,g,l,u,A,bl,bu,xinit)}. For Problem~8 the variable
bounds are infinite and the linear constraints are empty,
\begin{equation}
 l=-\infty\,\mathbf{1}_{2N},\quad u=+\infty\,\mathbf{1}_{2N},\quad
 A=0_{0\times2N},
 \label{eq:p8-qp-unconstrained}
\end{equation}
so the solver returns \eqref{eq:p8-solution}. Problems~9--10 reuse the same
call with finite bounds and a nonempty constraint matrix.

\section{Steady-State Target and Offset}
\label{sec:p8-target}

For a constant reference and disturbance estimate the target satisfies
\begin{equation}
 x_{ss}=Ax_{ss}+Bu_{ss}+G\hat w,
 \qquad
 r=Cx_{ss}.
 \label{eq:p8-ss-equations}
\end{equation}
Eliminating $x_{ss}$ gives
\begin{equation}
 Mu_{ss}=r-C(I-A)^{-1}G\hat w,
 \qquad
 M=C(I-A)^{-1}B,
 \qquad
 x_{ss}=(I-A)^{-1}\left(Bu_{ss}+G\hat w\right).
 \label{eq:p8-target-u}
\end{equation}
For the mass model $M$ takes the value
\begin{equation}
 M=\begin{bmatrix}0.1178&0.0589\\ 0.0948&0.1658\end{bmatrix}
 \ \mathrm{cm/(cm^3/s)}
 \label{eq:p8-M-numeric}
\end{equation}
and is nonsingular, so the target is unique. This target is what makes the design
offset-free: with the six-state estimator $\hat w_k$ enters through $G\hat w$, so
the target follows a step in the inflow and the output settles at the reference.
With the four-state estimator there is no $\hat w_k$, the target is computed for
$w=0$, and no state of the model can represent a persistent inflow, so an offset
remains. The stacked input target $\bar U_k=\mathbf{1}_N\otimes u_{ss}$ enters
\eqref{eq:p8-g} and is recomputed at every sample.

\section{Estimators Used by the MPC}
\label{sec:p8-estimators}

Both controllers use the stationary Kalman filters of Problem~6:
\texttt{mpc\_estimators} takes the steady-state gain from
\texttt{problem06Filter} and passes the gain, the disturbance dimension $n_w$ and
the prediction mode to the MPC, and \texttt{mpc\_estimate} advances the recursion
one sample. The Riccati equation and the augmented model are those of Problem~6
and are not repeated here.

For $n_w=0$ the filter is the four-state design, with prediction
$\hat x_{k|k-1}=A\hat x_{k-1|k-1}+Bu_{k-1}$, innovation
$e_k=y_k-C\hat x_{k|k-1}$ and update $\hat x_{k|k}=\hat x_{k|k-1}+K_{fx}e_k$,
where
\begin{equation}
 K_{fx}=\mathrm{diag}(19.50,20.19,10.34,12.08)
 \label{eq:p8-kf4-gain}
\end{equation}
and $\hat w_k$ is empty. For $n_w=2$ the state is augmented with the bias states
$w_k=[\Delta F_{3,k},\Delta F_{4,k}]^T$, which follow a random walk with
increment standard deviation $\sigma_B=2$~cm$^3$/s, and the recursion is
\begin{align}
 \hat x_{k|k-1}&=A\hat x_{k-1|k-1}+Bu_{k-1}+E\hat w_{k-1},
 & \hat w_{k|k-1}&=\hat w_{k-1},
 \label{eq:p8-kf6-pred}\\
 e_k&=y_k-C\hat x_{k|k-1},
 \label{eq:p8-kf6-innov}\\
 \hat x_{k|k}&=\hat x_{k|k-1}+K_{fx}e_k,
 & \hat w_{k|k}&=\hat w_{k-1}+K_{fw}e_k,
 \label{eq:p8-kf6-update}
\end{align}
with
\begin{equation}
 K_{fx}=\begin{bmatrix}80.49&0\\ 0&80.17\\ 84.64&0\\ 0&100.34\end{bmatrix},
 \qquad
 K_{fw}=\begin{bmatrix}3.55&0\\ 0&3.55\end{bmatrix}.
 \label{eq:p8-kf6-gains}
\end{equation}
The larger state gains and the nonzero bias gains make the augmented filter react
much more strongly to output errors, and its bias estimate enters the prediction
as a persistent input through $G=E$.

\section{Implementation}
\label{sec:p8-implementation}

The entry script \texttt{run\_mpc} sets the configuration of
Table~\ref{tab:p8-design} and calls \texttt{problem08Simulate}, which builds the
model, both estimators and both controllers, runs them on one noise realization,
computes the metrics and writes the MAT, PNG and PDF artifacts with
\texttt{saveExperiment}. Table~\ref{tab:p8-functions} maps the functions used.

\begin{table}[htbp]
 \centering
 \caption{Function map. All functions are in \texttt{code/problem08} except
 \texttt{problem06Filter}, the Problem~7 \texttt{qpsolver} and the shared
 export helpers.}
 \label{tab:p8-functions}
 \small
 \begin{tabular}{@{}p{0.27\linewidth}p{0.67\linewidth}@{}}
 \hline
 Function & Role and interface\\
 \hline
 \texttt{mpc\_constructors}
 & Builds $\Gamma$, $\Phi_x$, $\Phi_w$, the stacked weights, $\Lambda$,
   $I_0$, $H$ and $M_z,M_u,M_{\Delta u}$; checks the rank of $\Gamma$ and the
   definiteness of $H$; stores the QP bounds and constraint matrix:\newline
   \texttt{mpc = mpc\_constructors(}\newline
   \texttt{model, config, estimator)}.\\
 \texttt{mpc\_step}
 & One sample: estimates $(\hat x_k,\hat w_k)$, forms $b_k$, $c_k$, $g_k$,
   calls \texttt{qpsolver}, applies the first move and returns
   $U_k$, $Z_k$, $b_k$, $c_k$, $g_k$ and the innovation in \texttt{info}:\newline
   \texttt{[u,xhat,what,info] =}\newline
   \texttt{mpc\_step(mpc, xhat, what,}\newline
   \texttt{u\_prev, y, R, Ubar)}.\\
 \texttt{mpc\_predict}
 & Rolls out \eqref{eq:p8-rollout} over the horizon and verifies
   $Z=b+\Gamma U$ against the recursion:\newline
   \texttt{prediction = mpc\_predict(}\newline
   \texttt{mpc, xhat, what, U)}.\\
 \texttt{mpc\_target}
 & Solves \eqref{eq:p8-target-u} for the current reference and $\hat w_k$:\newline
   \texttt{[xss,uss] = mpc\_target(}\newline
   \texttt{mpc, r, what)}.\\
 \texttt{mpc\_estimators}, \newline \texttt{mpc\_estimate}
 & Builds the estimator for \texttt{'unbiased'} ($n_w=0$) or
   \texttt{'biased'} ($n_w=2$, persistent) from \texttt{problem06Filter}, and
   runs one prediction, innovation and update step of
   \eqref{eq:p8-kf6-pred}--\eqref{eq:p8-kf6-update}:\newline
   \texttt{estimator = mpc\_estimators(}\newline
   \texttt{kind, model, mode)};\newline
   \texttt{[xhat,what,e] =}\newline
   \texttt{mpc\_estimate(mpc, xprev,}\newline
   \texttt{wprev, uprev, y)}.\\
 \texttt{problem08Models}, \newline \texttt{problem08Simulate}
 & Builds the Problem~5 mass model and the identified Problem~4 model, and runs
   both controllers and the metrics:\newline
   \texttt{[result,figures] =}\newline
   \texttt{problem08Simulate(config)}.\\
 \hline
 \end{tabular}
\end{table}

The constructor is model-agnostic. Besides the mass model it is applied to the
six-state zero-order-hold realization of the identified transfer function from
Problem~4, discretized at $T_s=5$~s with zero feedthrough; that design has
$N=20$ and positive-definite $H$, so the framework carries over to an identified
model without changes. It is not simulated, because the Problem~6 estimators are
defined for the mass model, and the comparison below therefore uses the mass
model only.

\section{Results}
\label{sec:p8-results}

The reported run uses seed $8001$ and the weights of
Table~\ref{tab:p8-design}. The reference is a step of $[1,0.5]^T$~cm at
$t=200$~s; the unmeasured inflow steps are $\Delta F_3=20$~cm$^3$/s at $t=300$~s
and $\Delta F_4=15$~cm$^3$/s at $t=500$~s, with the noise of
\eqref{eq:p8-noise}. Both controllers run on the linear plant
\eqref{eq:p8-model-state}--\eqref{eq:p8-model-out} with the same noise
realization, so the comparison is paired; no constraints are imposed. The
metrics are collected in Table~\ref{tab:p8-metrics}.

\begin{table}[htbp]
 \centering
 \caption{Closed-loop metrics for the linear plant with seed $8001$.
 RMSE and maximum error are over $t\ge200$~s; the offset is the mean over
 the last $100$~s. Effort sums are over the full run, in
 $(\mathrm{cm^3/s})^2$.}
 \label{tab:p8-metrics}
 \small
 \begin{tabular}{@{}lrrrrrrr@{}}
 \hline
 Controller & \multicolumn{2}{c}{RMSE (cm)} & \multicolumn{2}{c}{Offset (cm)}
 & $\max|e|$ (cm) & $\sum u_k^2$ & $\sum(\Delta u_k)^2$\\
 & $h_1$ & $h_2$ & $h_1$ & $h_2$ & & & \\
 \hline
 unbiased 4-state & 1.965 & 1.197 & $+2.235$ & $+1.756$ & 2.321
  & $12\,683$ & $200.1$\\
 biased 6-state & 0.334 & 0.281 & $-0.221$ & $-0.101$ & 1.079
  & $131\,513.6$ & $1\,906.6$\\
 \hline
 \end{tabular}
\end{table}

\paragraph{Offset and tracking.}
The four-state controller leaves a final offset of $+2.235$~cm in $h_1$ and
$+1.756$~cm in $h_2$, against $-0.221$ and $-0.101$~cm for the six-state
controller, and its RMSE falls from $1.965$/$1.197$~cm to $0.334$/$0.281$~cm.
Figure~\ref{fig:p8-tracking} shows the same contrast in the time domain. A
near-deterministic check with $\sigma_F=\sigma_V=10^{-3}$ gives offsets
$[3.867,3.437]^T$~cm and $[-0.0006,-0.0004]^T$~cm, so the four-state offset is
structural, not a noise artifact: its target \eqref{eq:p8-target-u} is computed
for the wrong disturbance, whereas the six-state target includes the estimated
bias and removes the offset.

\paragraph{Effort and disturbance estimates.}
The six-state controller uses about ten times the input energy, $131\,513.6$
versus $12\,683$ in $\sum u_k^2$ (ratio $10.4$) and $1\,906.6$ versus $200.1$ in
$\sum(\Delta u_k)^2$ (ratio $9.5$). Figure~\ref{fig:p8-pumps} shows the cause:
apart from the transient after the reference step the four-state pumps stay near
the nominal $300$~cm$^3$/s, while the six-state pumps follow the changing target,
with pump~1 falling towards $260$~cm$^3$/s after the first disturbance step and
pump~2 rising to about $330$~cm$^3$/s around the second. This is the price of the
integral action. The final disturbance estimates in
Figure~\ref{fig:p8-disturbance} are $\hat F_3=22.63$ and
$\hat F_4=17.92$~cm$^3$/s against the true $20$ and $15$, i.e.\ errors of
$+13\%$ and $+19\%$, with a lag of tens of seconds; the residual error is of the
order of the random-walk increment used for design. It shifts the input target
but does not by itself produce an output offset, because the offset-free property
comes from the integrator and the target calculation.

\paragraph{Prediction.}
Figure~\ref{fig:p8-prediction} shows the six-state prediction at $t=700$~s,
after both disturbance steps. The $20$-step prediction covers $100$~s and
converges to the $[1,0.5]^T$~cm reference. The realized trajectory departs from
it by a few tenths of a centimetre because the plant sees the actual noisy
inflow, whereas the prediction holds the current bias estimate constant: it is
the model-consistent trajectory used by the optimizer, not a forecast of the
noise.

\begingroup
\newcommand{\pEightFigure}[1]{%
 \includegraphics[width=\linewidth]{tex/chapters/results/problem08/#1.pdf}}
\begin{figure}[htbp]
 \centering\pEightFigure{unconstrained_mpc_01}
 \caption{Closed-loop levels for the four-state MPC (left column) and the
 six-state MPC (right column). Rows are $h_1$ and $h_2$; solid black is the
 true level, dashed blue the reference, and grey points the noisy measurement.}
 \label{fig:p8-tracking}
\end{figure}
\begin{figure}[htbp]
 \centering\pEightFigure{unconstrained_mpc_02}
 \caption{$F_3$ and $F_4$ estimates from the six-state MPC (dashed) compared
 with the true disturbance steps (solid).}
 \label{fig:p8-disturbance}
\end{figure}
\begin{figure}[htbp]
 \centering\pEightFigure{unconstrained_mpc_03}
 \caption{Physical pump flows $u_0+u$ for the two controllers. The six-state
 controller moves the pumps much more because its target is recomputed from the
 disturbance estimate.}
 \label{fig:p8-pumps}
\end{figure}
\begin{figure}[htbp]
 \centering\pEightFigure{unconstrained_mpc_04}
 \caption{$20$-step ($100$~s) horizon prediction of the six-state MPC at
 $t=700$~s (dashed with markers) and the realized levels (solid).}
 \label{fig:p8-prediction}
\end{figure}
\endgroup

\section{Discussion}
\label{sec:p8-discussion}

All gains are stationary, so the estimator is linear time-invariant and $H$ can
be precomputed; the price is that the controller treats $\hat w_k$ as certain
and that $\sigma_B=2$~cm$^3$/s must be a single scalar rather than a
state-dependent weight. This value sets the bandwidth of the disturbance
estimator: larger values follow the steps faster but give a noisier estimate and
larger input moves, smaller values give smoother signals and a longer offset
transient. The residual estimate errors and the residual offset reflect this
trade-off, and the value is a tuning choice rather than the result of an
optimization. The comparison uses one noise realization, so the metrics are
paired but not averaged over seeds. Finally, the experiment is on the linear
plant by design: it isolates the estimator and the target calculation, and input
and output constraints are the subject of Problems~9--10.
